%Zeigerdiagramme


\newvideofile{Komplx}{Komplexe Zahlen}

\begin{frame}
    \fta{Fundamentals of Complex Numbers}

     
    \s{
        The phasor diagrams used in alternating current technology provide a quick overview of the magnitude and
        direction of the voltage and current. The underlying concepts of the complex number plane and
        the structure of phasor diagrams are explained in more detail below and illustrated with an example.
    }



    \begin{Learning objectives}{Complex numbers}
        The Students can
        \begin{itemize}
            \item handle numbers in the complex plane.
            \item Plot phasor diagrams of complex numbers.
            \item calculate complex numbers.
        \end{itemize}
    \end{Learning objectives}

    \speech{KomplxLrnz}{1}{Das Themengebiet der komplexen Zahlen ist notwendig, 
    um tiefer in die Gegebenheiten der Wechselstromlehre und ihrer periodischen Größen einzutauchen.
    Hierzu sollen in diesem Kapitel die Grundlagen der Komplexen Zahlen vermittelt werden.
    Anschließend soll mit Zahlen in der komplexen Ebene umgegangen,
    Zeigerdiagramme von komplexen Zahlen dargestellt 
    und komplexe Zahlen berechnet werden können.} 
\end{frame}

\begin{frame}
    \ftb{Complex number plane}

    \s{
        In order to understand the structure of pointer diagrams, a basic understanding of the complex 
        number system is necessary. For elementary arithmetic operations, natural numbers including zero and 
        rational numbers are sufficient. Rational numbers can be represented as finite or periodic decimal numbers.
        Irrational numbers, on the other hand, can be represented as decimal numbers that have an infinite number of
        digits and are not periodic. Real numbers are composed of rational numbers
        and irrational numbers. However, it is not possible to take the square root of a negative number 
        in the real number system. For this purpose, the complex number system is introduced. In the complex 
        number system, the real number space is extended by the imaginary unit j. Thus, in the complex
        number system, taking the square root of -1 results in the imaginary unit j.
    }

    \b{
        In the complex number plane, the imaginary unit j is introduced in order to perform calculations with complex numbers.
    }


    
    \begin{eq}
        \mathrm{j}=\sqrt{-1} 
    \end{eq}

    \onslide<2->{

        Squaring the imaginary unit again gives -1. 

    \begin{eq}
        \mathrm{j}^2=-1
    \end{eq}}

    \speech{KomplxImag}{1}{In der Mathematik besteht das Problem, das die Wurzel aus einer negativen Zahl im reellen Zahlenraum nicht weiter definiert ist.
    Hier wird in der Elektrotechnik die Imaginäre Einheit jot eingeführt, welche in anderen Themenbereichen auch I genannt wird.} 
    \speech{KomplxImag}{2}{Die Definition der angegebenen Gleichung, das jot zum Quadrat gleich minus eins ergibt, ermöglicht das Radizieren von negativen Zahlen.}
\end{frame}


\begin{frame}
    \ftx{Representation of complex numbers}

    \s{
        A complex number \underline{Z} describes a location in the complex plane. To uniquely define a location in a two-dimensional
        coordinate system, two coordinates are required. 
        The two coordinates for describing a complex number \underline{Z} are described in the complex plane as the real part and the imaginary part
        (see Equation \ref{GleichungKomplexeZahlen}). Complex numbers are usually indicated by an underscore,
        whereby the real part and the imaginary part represent real numbers. 
    }

    \b{
        Cartesian coordinates with the real part and the product of the imaginary unit and the imaginary part:
    }

    \begin{eq}
        \underline{Z}=Real part+\mathrm{j} \cdot Imaginary\ part \label{GleichungKomplexeZahlen}
    \end{eq}

    \onslide<2->{
    \b{
        The real part is abbreviated with $\Re$ and the imaginary part with $\Im$:
    }

    \begin{eq}
        \underline{Z}=\Re(\underline{Z})+\mathrm{j} \cdot \Im(\underline{Z})  
    \end{eq}}

    \speech{KomplxDarst}{1}{Der skalare reelle Zahlenbereich wird um den Bereich der komplexen Zahlen erweitert und führt so zur komplexen Zahlenebene.
    Hier kann die komplexe Zahl beispielsweise mit kartesischen Koordinaten beschrieben werden. Die komplexe Zahl Zett unterteilt sich in den Realteil,
    gepaart mit der imaginären Einheit und dem Imaginärteil. Die Komplexen Zahlen werden mit einem Unterstrich gekennzeichnet.}
    \speech{KomplxDarst}{2}{Abgekürzt wird der Realteil Er Ee. Gelesen: Realteil von Zett. Der Imaginärteil wird mit I emm abgekürzt.}
\end{frame}


\begin{frame}
    \renewcommand{\figurename}{Figure}
    \ftx{Representation of complex numbers}

    \s{
        In the complex plane, the real part is plotted on the abscissa and the imaginary part on the ordinate. 
        The abbreviations Re = real part and Im = imaginary part are used. The coordinate system of a complex 
        plane is explained in Figure \ref{BildKomplexeEbene}. 
        The location of the complex number can be represented by 
        direction arrows, also known as vectors or pointers. The representation
        of the complex number in Cartesian coordinates is achieved by decomposing it into the real part of the complex number 
        Re(\underline{Z}) and the imaginary part of the complex number Im(\underline{Z}) combined with the imaginary unit j.
        Complex numbers can also be represented in polar coordinates. This is done by the magnitude $|\underline{Z}|$
        of the complex number \underline{Z} and by the angle $\varphi$ that the pointer of the complex number encloses with the real axis. 
    }

    \b{
        Use of the complex layer:
        \begin{itemize}
            \item<1-> The real part $\Re$ is plotted on the vertical axis and the imaginary part $\Im$ on the horizontal axis.
            \item<2-> Cartesian coordinates: $\Re(\underline{Z}) + \mathrm{j} \cdot \Im(\underline{Z})$
            \item<3-> Polar coordinates: $|\underline{Z}| \cdot e^{j \cdot \varphi_\mathrm{Z}}  $
        \end{itemize}
    }

    \fu{
        \input{Tikz/src/Komplexe_Zahlen/Zeigerdiagramm_Komplexe_Zahlen}
    }{{\bf Pointer diagram of a complex number.} A complex number $\underline{Z}$, which can be represented by $Re(\underline{Z})$ and $Im(\underline{Z})$
    in Cartesian coordinates and by the magnitude of the complex number $|\underline{Z}|$ and the corresponding angle
    $\varphi$ in polar coordinates. \label{BildKomplexeEbene}}

    \speech{KomplxDarstEbn}{1}{Folgend wird die komplexe Ebene dargestellt. Hier wird der Realteil auf der Abszisse und der Imaginärteil auf der Ordinate aufgetragen.}
    \speech{KomplxDarstEbn}{2}{Eine komplexe Zahl Zett kann in der komplexen Ebene dargestellt werden. Nach der zuvor getroffenen Definition, 
    kann die komplexe Zahl in den Realteil von Zett und den Imaginärteil von Zett unterteilt werden.}
    \speech{KomplxDarstEbn}{3}{Außerdem kann eine komplexe Zahl auch in Polarkoordinaten dargestellt werden.
    Hierbei wird der Betrag der komplexen Zahl und der Winkel Vieh zur Abszisse als Definition verwendet.}
\end{frame}


\begin{frame}
    \ftx{Coordinate transformation}

    \s{
        The two forms of representation can be transformed into each other. Euler's formula is helpful here.
        Euler's formula shows that the ordinate value and the abscissa value of the unit circle
        can be calculated using the trigonometric functions cosine and sine. 
        Euler's formula is represented in equation \ref{GleichungEuler}.

        \begin{eq}
            \mathrm{e}^{\mathrm{j} \cdot \chi} = \cos(\chi) + \mathrm{j} \cdot \sin(\chi)    \label{GleichungEuler}
        \end{eq}
    
        Using Euler's formula, the polar coordinates
        can be transformed into Cartesian coordinates by adding the absolute value of the complex number. This is explained by equation \ref{GleichungKartesisch}.
     
        \begin{eq}
            \underline{Z} = |\underline{Z}| \cdot \cos(\varphi) + \mathrm{j} \cdot |\underline{Z}| \cdot \sin(\varphi) = \Re(\underline{Z}) + \mathrm{j} \cdot \Im(\underline{Z}) \label{GleichungKartesisch}
        \end{eq}
    
        Using Pythagoras' theorem, the real and imaginary parts of the complex number can be used to determine the magnitude of the complex number
        for the polar coordinates. The corresponding angle is obtained from the arctan with the ratio of the imaginary part
        to the real part in the argument. With this information, a complex number such as in equation \ref{GleichungPolar} 
        can be transformed into polar coordinates.
 
    
        \begin{eq}
            \underline{Z} = \sqrt{(\Re(\underline{Z}))^2+(\Im(\underline{Z}))^2} \cdot \mathrm{e}^{\mathrm{j} \cdot \arctan(\frac{\Im(\underline{Z})}{\Re(\underline{Z})}) } = |\underline{Z}| \cdot \mathrm{e}^{\mathrm{j} \cdot \varphi} \label{GleichungPolar}
        \end{eq}
    }

    \b{
        \begin{itemize}
            \item Euler's formula: 
            \begin{eq}
                \mathrm{e}^{\mathrm{j} \cdot \chi} = \cos(\chi) + \mathrm{j} \cdot \sin(\chi)  
            \end{eq}
            \item Cartesian coordinates:
            \begin{eq}
                \underline{Z} = |\underline{Z}| \cdot \cos(\varphi) + \mathrm{j} \cdot |\underline{Z}| \cdot \sin(\varphi) = \Re(\underline{Z}) + \mathrm{j} \cdot \Im(\underline{Z}) 
            \end{eq}
            \item Polar coordinates:
            \begin{eq}
                \underline{Z} = \sqrt{(\Re(\underline{Z}))^2+(\Im(\underline{Z}))^2} \cdot \mathrm{e}^{\mathrm{j} \cdot \arctan(\frac{\Im(\underline{Z})}{\Re(\underline{Z})}) } = |\underline{Z}| \cdot e^{j \cdot \varphi} 
            \end{eq}
        \end{itemize}
    }

    \speech{KomplxKoordtr}{1}{Über die Eulersche Formel können Winkelangaben, in die Bestandteile der Kartesischen Koordinaten aufgeteilt werden.
    Um die komplexe Zahl aus den Polarkoordinaten in die kartesischen Koordinaten zu überführen, 
    wird über die Eulersche Formel mit dem Betrag und dem Winkel der komplexen Zahl der Realteil und der Imaghinärteil bestimmt.
    Andersherum lässt sichd er Betrag der komplexen Zahl über den Satz des Pütagoras bestimmen. Der Winkel ergibt sich aus der Arcustangensfunktion.}
\end{frame}


\begin{frame}
    \ftx{Representations of complex numbers}

    \begin{Key point}{Representations of complex numbers}
        Cartesian representation:
        \begin{eq}
            \underline{Z} = \Re(\underline{Z}) + j \cdot \Im(\underline{Z}) \nonumber
        \end{eq}
        Polar form:
        \begin{eq}
            \underline{Z} = |\underline{Z}| \cdot e^{j \cdot \varphi}   \nonumber
        \end{eq}
        Trigonometric representation:
        \begin{eq}
            \underline{Z} = |\underline{Z}| (\cos(\varphi) + j \cdot \sin(\varphi))   \nonumber
        \end{eq}
    \end{Key point}

    \speech{KomplxMrk1}{1}{Komplexe Zahlen lassen sich über den Realteil und Imaginärteil in Kartesischen Koordinaten,
    über die Polarform mit dem Betrag und dem zugehörigen Winkel,
    sowie mit den Trigonometrischen Funktionen darstellen und ineinander umrechnen.}
\end{frame}

\begin{frame}
    \ftx{Basic arithmetic and operations}

    \s{
        Conjugation is one of the operations on complex numbers. When conjugating a complex number, the 
        j is replaced by -j (negation of the imaginary part). As shown in Equation \ref{GleichungKonj}, conjugation is indicated by an asterisk.
    }

    \b{
        Conjugation:
    }

    \begin{eq}
        \underline{Z}^* = (\Re(\underline{Z}) + \mathrm{j} \cdot \Im(\underline{Z}))^*= \Re(\underline{Z}) - \mathrm{j} \cdot \Im(\underline{Z})   \label{GleichungKonj}
    \end{eq}

    \s{
        The magnitude of a complex number can be determined by multiplying the complex number by its conjugate
        (see Equation \ref{GleichungBetrag}). Instead of the magnitude, which contains a root, the square of the magnitude is also used.
        It also represents a measure of the distance of the number from the origin, but is easier to calculate.
    }

    \onslide<2->{
    \b{
        Amount and amount squared
    }

    \begin{eq}
        |\underline{Z}| = \sqrt{\underline{Z} \cdot \underline{Z}^*} \rightarrow |\underline{Z}|^2 = \underline{Z} \cdot \underline{Z}^*        \label{GleichungBetrag}
    \end{eq}}

    \speech{KomplxOperat}{1}{Bei der sogennanten Konjugiert-Komplexen Darstellung einer Zahl handelt es sich um eine komplexe Zahl, 
    bei welcher das Vorzeichen des Imaginärteiles negiert wird. So wird beispielsweise aus dem positiven Imaginärteil ein negativer.}
    \speech{KomplxOperat}{2}{Das Betragsquadrat quadriert den Betrag einer Komplexen Zahl. 
    Es wird ebenso wie der normale Betrag einer komplexen Zahl als Maß für den Abstand zum Ursprung verwendet. }
\end{frame}

\begin{frame}
    \ftx{Basic arithmetic and operations}
    \s{
        When adding and subtracting complex numbers, it is advisable to first convert them into Cartesian coordinates. 
        This allows the real part and the imaginary part of the complex number
        to be calculated separately.
    }
    \b{
        Addition and subtraction:
    }
    \begin{eqa}
        \underline{Z}_1+\underline{Z}_2 = \Re(\underline{Z}_1) + \Re(\underline{Z}_2) +j \cdot (\Im(\underline{Z}_1) + \Im(\underline{Z}_2)) \\
        \underline{Z}_1-\underline{Z}_2 = \Re(\underline{Z}_1) - \Re(\underline{Z}_2) -j \cdot (\Im(\underline{Z}_1) - \Im(\underline{Z}_2))
    \end{eqa}

    \s{
        The representation in polar coordinates is recommended for calculating complex numbers by multiplication
        and division. The multiplication of complex numbers is composed of the product of the magnitudes and the sum of the respective
        angles (see Equation \ref{GleichungMultiplikation}). In division, the magnitudes are divided and the angles are subtracted from each other (see Equation \ref{GleichungDivision}). 
    }

    \onslide<2->{
    \b{
        Multiplication and division:
    }

    \begin{eq}
        \underline{Z}_1 \cdot \underline{Z}_2 = |\underline{Z}_1| \cdot |\underline{Z}_2| \cdot e^{j \cdot (\varphi_1+\varphi_2)} \label{GleichungMultiplikation}
    \end{eq}

    \begin{eq}
        \frac{\underline{Z}_1}{\underline{Z}_2} = \frac{|\underline{Z}_1|}{|\underline{Z}_2|} \cdot e^{j \cdot (\varphi_1-\varphi_2)} \label{GleichungDivision}
    \end{eq}}

    \speech{KomplxGrundrech}{1}{Grundsätzlich wird bei verschiedenen komplexen Zahlen gerechnet wie mit unterschiedlichen Variablen. 
    Bei der Addition werden die Realteile und imaginärteile jeweils seperat addiert.
    Ebenso verhält es sich bei der Subtraktion, auch hier werden die Realteile und Imaginärteile jeweils für sich voneinander subtrahiert.}
    \speech{KomplxGrundrech}{2}{Bei der Multiplikation wird der Betrag der Komplexen Zahlen miteinander Multipliziert und die Winkel im Exponenten addiert.
    Die Division von komplexen Zahlen erfolgt durch die Division der Beträge und der Subtraktion der Winkel.}
\end{frame}


\begin{frame}
    \ftx{Basic arithmetic operations with complex numbers}

    \begin{Key point}{Basic arithmetic operations with complex numbers}
        For {\bf addition and subtraction}, the {\bf Cartesian representation} is generally used for complex numbers.
        For {\bf multiplication and division} of complex numbers, the {\bf polar form} is chosen.
    \end{Key point}

    \speech{KomplxMrk2}{1}{Für die Addition und die Subtraktion wird in der Regel die kartesische Darstellungsform für komplexen Zahlen verwendet.
    Bei der Multiplikation und der Division von komplexen Zahlen wird die Polarform zur Berechnung gewählt.}
\end{frame}


\begin{frame}
    \ftx{Basic arithmetic and operations}
    \s{
        Since the division is initially only defined for real numbers other than zero, the complex number in the denominator is extended so that 
        it becomes real.
    }
    \b{
        Division by conjugate complex extension:
    }
    \begin{eq}
        \frac{\underline{Z}_1}{\underline{Z}_2} = \frac{\underline{Z}_1}{\underline{Z}_2} \cdot \frac{\underline{Z}_2^*}{\underline{Z}_2^*} = \frac{\underline{Z}_1 \cdot \underline{Z}_2^*}{|\underline{Z}_2|^*} 
    \end{eq}

    \s{
        Adding a complex number to its conjugate cancels out its imaginary part (see equation \ref{GleichungRealteil}). 
        Similarly, subtracting the real part cancels it out (see equation \ref{GleichungImaginärteil}).    
    }

    \onslide<2->{
    \b{
        Calculate the real part and imaginary part:
    }
    \begin{eq}
        \Re(\underline{Z}) = \frac{\underline{Z}+\underline{Z}^*}{2} \label{GleichungRealteil}
    \end{eq}
    \begin{eq}    
        \Im(\underline{Z}) = \frac{\underline{Z}-\underline{Z}^*}{2j} \label{GleichungImaginärteil}
    \end{eq}}

    \speech{KomplxWeitere}{1}{Die Division von komplexen Zahlen kann auch durch die konjugiert-komplexe Erweiterung erfolgen.} 
    \speech{KomplxWeitere}{2}{Durch die Addition einer komplexen Zahl mit ihrer konjugiert komplexen hebt sich ihr Imaginärteil auf. 
    Entsprechend verschwindet durch die Subtraktion der Realteil.
    Je nach Anwendungsfall können mit diesen Operationen die Rechnungen vereinfacht werden.}     
\end{frame}


\begin{frame}
    \renewcommand{\figurename}{Figure}
    \ftb{Graphical representation and calculations with complex numbers}

    \s{
        Complex numbers can also be solved graphically using pointer diagrams. If the complex
        numbers are represented in Cartesian form as recommended, the components of the real parts and the 
        imaginary parts can be read directly and entered into a coordinate system. The complex numbers 
        are entered, for example, in Figure \ref{BildKomplexRechnung}. When adding complex numbers, either 
        $Z_\mathrm{1}$ is set to $Z_\mathrm{2}$ or vice versa. This is because the commutative law still applies to addition. Here, 
        the dashed arrows indicate the parallel shifts of $Z_\mathrm{1}$ and $Z_\mathrm{2}$. Both shifts end at 
        the same coordinate. Here, the real part and the imaginary part can be read separately from the sum of the two complex 
        numbers.
    }

    \b{
        Arithmetic operations in the pointer diagram - Addition:

        \begin{itemize}
            \item[1.] <2-> Parallel translation $Z_\mathrm{2}$ to the tip of $Z_\mathrm{1}$
            \item[2.] <3-> Read off the real part and imaginary part
        \end{itemize}
    }

    \fu{
        \input{Tikz/src/Komplexe_Zahlen/Addition_Komplexer_Zahlen}
    }{{\bf Addition of complex numbers.} Graphical solution for adding two complex numbers in a phasor diagram. \label{BildKomplexRechnung}}

    \speech{KomplxGraphAdd}{1}{Gezeigt wird die Komplexe Ebene mit der Aufteilung in Realteil und Imaghinärteil.
    Außerdem werden die Zahlen Zett eins in rot und Zett zwei in blau dargestellt.
    Hier soll nun die Addition der beiden komplexen Zahlen grafisch gelöst werden.}
    \speech{KomplxGraphAdd}{2}{Dafür wird eine der Zahlen parallel an der anderen komplexen Zahl verschoben. 
    Dabei ist es nach dem Kommutativgesetz unerheblich, welche komplexe Zahl womit addiert wird.}
    \speech{KomplxGraphAdd}{3}{Nun können der Realteil und der Imaginärteil aus der Summe der beiden komplexen Zahlen abgelesen und notiert werden.}
\end{frame}


\begin{frame}
    \ftx{Pointer diagrams - Subtraction}

    \s{
        The graphical solution for subtracting complex numbers is similar to that for addition.
        It should be noted that, as in other number systems, the commutative law does not apply when subtracting two complex numbers. 
        If the complex number $Z_\mathrm{2}$ is subtracted from $Z_\mathrm{1}$, the direction
        of $Z_\mathrm{1}$ changes. Here, the parallel shift occurs again, but only for $Z_\mathrm{1}$. The resulting
        new vector can then be read again as the real part and the imaginary part.
    }

    \b{
        Arithmetic operations in the pointer diagram - Subtraction:

        \begin{itemize}
            \item[1.] <1-> Pointer reversal for negation
            \item[2.] <2-> Move pointer $Z_\mathrm{1}$: Base point of $Z_\mathrm{1}$ to tip of $Z_\mathrm{2}$
            \item[3.] <3-> Read off the real part and imaginary part
        \end{itemize}
    }

    \fu{
        \input{Tikz/src/Komplexe_Zahlen/Subtraktion_Komplexer_Zahlen}
    }{{\bf Subtraction of complex numbers.} Graphical solution for subtracting two complex numbers in a phasor diagram.}

    \speech{KomplxGraphSub}{1}{Ähnlich wie bei der Addition wird bei der Subtraktion von komplexen Zahlen vorgegangen.
    Wieder wird die komplexe Ebene mit dem Realteil und dem Imaginärteil dargestellt.
    Nun wird aber Zett eins von Zett zwei subtrahiert. 
    Hierfür wird die Richtung des Zeigers Zett eins umgekehrt. }
    \speech{KomplxGraphSub}{2}{Anschließend wird der Zeiger von Zett eins so verschoben,
    dass das Ende von Zett eins an der Spitze von Zett zwei anliegt.}
    \speech{KomplxGraphSub}{3}{Folgend werden wieder der Realteil und der Imaginärteil aus der Differenz der beiden komplexen Zahlen abgelesen und notiert.}
\end{frame}

  
\begin{frame}
    \ftx{Pointer diagrams - Multiplication}

    \s{
        The complex numbers $\underline{Z}_\mathrm{1}$ and $\underline{Z}_\mathrm{2}$ shown in Figure \ref{BildMultiplikation} are represented for 
        multiplication in polar coordinates. Using the previously described calculation rules for multiplication 
        of complex numbers, the values for the result $\underline{Z}_\mathrm{3}$ can be determined. The multiplication
        of the pointer lengths gives the length of the product. The sum of the two complex numbers gives the angle of the new
        complex number.
    }

    \b{
        Arithmetic operations in the pointer diagram - Multiplication: (Drawing not to scale!)

        \begin{itemize}
            \item[1.] <2-> Winkel $\varphi_\mathrm{1}+\varphi_\mathrm{2}$ berechnen und einzeichnen
            \item[2.] <3-> Calculate and enter the length $Z_1 \cdot Z_\mathrm{2}$
            \item[3.] <4-> Read off the real part and imaginary part
        \end{itemize}
    }
    
    \fu{
        \input{Tikz/src/Komplexe_Zahlen/Multiplikation_Komplexer_Zahlen}
    }{{\bf Multiplication of complex numbers.} Graphical solution for multiplying two complex numbers in polar coordinates in a pointer diagram. \label{BildMultiplikation}}

    \speech{KomplxGraphMulti}{1}{Bei der grafischen Multiplikation ist ein wenig Rechnen notwendig.
    Hier werden die komplexen Zahln Zett eins und Zett zwei in der komplexen Ebene in ihrer Polarform dargestellt.}
    \speech{KomplxGraphMulti}{2}{Um das Produkt aus den beiden komplexen Zahlen zu bestimmen,
    wird im ersten Schritt die Summe aus den beiden Winkel der komplexen Zahlen gebildet und notiert.}
    \speech{KomplxGraphMulti}{3}{Im zweiten Schritt werden die beiden Beträgen der komplexen Zahlen miteinander multipliziert und eingetragen.}
    \speech{KomplxGraphMulti}{4}{Abschließend kann das Ergebniss nach der Berechnung auch grafisch veranschaulicht werden.}
\end{frame} 
\newpage